\documentclass[11pt,a4paper]{article}
\usepackage[utf8]{inputenc}
\usepackage{amsmath,amssymb,amsthm}
\usepackage{listings}
\usepackage{xcolor}
\usepackage[margin=1in]{geometry}
\usepackage{hyperref}

\newtheorem{theorem}{Theorem}

\lstset{
  language=C++,
  basicstyle=\ttfamily\small,
  keywordstyle=\color{blue}\bfseries,
  commentstyle=\color{green!60!black},
  stringstyle=\color{red},
  numbers=left,
  numberstyle=\tiny\color{gray},
  breaklines=true,
  frame=single,
  tabsize=4,
  showstringspaces=false
}

\title{IOI 2005 -- Garden (Largest Empty Rectangle)}
\author{Solution}
\date{}

\begin{document}
\maketitle

\section{Problem Statement Summary}
Given an $R \times C$ grid with some cells blocked, find the area of the
largest axis-aligned rectangle consisting entirely of unblocked cells.

\section{Solution: Largest Rectangle in Histogram}

\subsection{Algorithm}
\begin{enumerate}
  \item For each cell $(i, j)$, compute $h[j]$ = the number of
        consecutive unblocked cells in column $j$ ending at row $i$
        (inclusive). Reset $h[j] = 0$ if $(i, j)$ is blocked.
  \item For each row, treat $h[0], h[1], \ldots, h[C{-}1]$ as a
        histogram and find the area of the largest rectangle in it.
  \item The answer is the maximum over all rows.
\end{enumerate}

\subsection{Stack-Based Histogram Algorithm}
For a histogram of $C$ bars:
\begin{enumerate}
  \item Maintain a stack of bar indices with non-decreasing heights.
  \item For each bar $j$ (and a sentinel bar of height 0 at position
        $C$): while the stack top has height $\ge h[j]$, pop it and
        compute the rectangle of that height, extending from the new
        stack top $+ 1$ to $j - 1$.
  \item Each bar is pushed and popped at most once, giving amortized
        $O(C)$.
\end{enumerate}

\section{C++ Implementation}

\begin{lstlisting}
#include <bits/stdc++.h>
using namespace std;

int largestRectInHistogram(vector<int>& h) {
    int n = h.size();
    stack<int> st;
    int maxArea = 0;

    for (int i = 0; i <= n; i++) {
        int cur = (i == n) ? 0 : h[i];
        while (!st.empty() && h[st.top()] >= cur) {
            int height = h[st.top()];
            st.pop();
            int width = st.empty() ? i : (i - st.top() - 1);
            maxArea = max(maxArea, height * width);
        }
        st.push(i);
    }
    return maxArea;
}

int main() {
    ios::sync_with_stdio(false);
    cin.tie(nullptr);

    int R, C;
    cin >> R >> C;

    vector<vector<int>> grid(R, vector<int>(C));
    for (int i = 0; i < R; i++)
        for (int j = 0; j < C; j++)
            cin >> grid[i][j];

    vector<int> h(C, 0);
    int ans = 0;

    for (int i = 0; i < R; i++) {
        for (int j = 0; j < C; j++) {
            if (grid[i][j] == 0) // free cell
                h[j]++;
            else
                h[j] = 0;
        }
        ans = max(ans, largestRectInHistogram(h));
    }

    cout << ans << "\n";
    return 0;
}
\end{lstlisting}

\section{Complexity Analysis}
\begin{itemize}
  \item \textbf{Time:} $O(R \times C)$. One $O(C)$ histogram pass per
        row.
  \item \textbf{Space:} $O(C)$ (height array and stack), or $O(R \times C)$
        if the full grid is stored.
\end{itemize}

This is optimal since reading the input alone requires $O(R \times C)$.

\paragraph{Note.} The specific IOI 2005 ``Garden'' problem may impose
additional constraints (e.g., the rectangle must contain at least $K$
specific points, or two rectangles must be found). The histogram
technique serves as the core building block, augmented by binary search
or two-pointer methods as needed.

\end{document}
